\begin{tikzpicture}[x={\dimexpr\linewidth/169\relax},y=1mm,
  font=\figurefont\footnotesize,text=NNInk,
  mc flow/.style={nn flow},
  mc channel/.style={fill=none,inner sep=.4pt,font=\figurefont\footnotesize},
  mc sum/.style={circle,draw=NNHidden,fill=white,line width=1.1pt,
    minimum size=6mm,inner sep=0pt}]
  % 同一输入的三个通道沿对角线错位叠放；窗口位于各平面的外露部分。
  \foreach \c/\xx/\yy/\col in {3/18/-4/NNOutput,2/9/-13/NNHidden,1/0/-22/NNInput}{
    \draw[draw=\col,fill=\col!7,line width=1.1pt]
      (\xx,\yy) rectangle ({\xx+24},{\yy+24});
    \foreach \g in {1,2,3}{
      \draw[draw=\col!45,line width=.5pt]
        ({\xx+6*\g},\yy)--({\xx+6*\g},{\yy+24})
        (\xx,{\yy+6*\g})--({\xx+24},{\yy+6*\g});
    }
    \draw[draw=\col,fill=\col!30,line width=1.05pt]
      ({\xx+12},{\yy+12}) rectangle ({\xx+24},{\yy+24});
    \draw[draw=\col,line width=.55pt]
      ({\xx+18},{\yy+12})--({\xx+18},{\yy+24})
      ({\xx+12},{\yy+18})--({\xx+24},{\yy+18});
    \coordinate (input\c) at ({\xx+24},{\yy+18});
  }
  \node[mc channel] at (7,-10) {$c=1$};
  \node[mc channel] at (16,7.5) {$c=2$};
  \node[mc channel] at (25,16.5) {$c=3$};
  \node[nn heading,anchor=north] at (21,-26) {$\bm X$};
  % 两个输出核独立展开，每组的三个二维切片竖排。
  \foreach \o/\base in {1/26,2/-26}{
    \foreach \c/\dy/\col in {3/12/NNOutput,2/0/NNHidden,1/-12/NNInput}{
      \pgfmathsetmacro{\ky}{\base+\dy}
      \draw[draw=\col,fill=\col!20,line width=1.1pt]
        (74,{\ky-4.5}) rectangle (83,{\ky+4.5});
      \draw[draw=\col,line width=.6pt]
        (78.5,{\ky-4.5})--(78.5,{\ky+4.5})
        (74,\ky)--(83,\ky);
      \coordinate (kernel\o\c-in) at (74,\ky);
      \coordinate (kernel\o\c-out) at (83,\ky);
    }
    \node[mc sum] (sum\o) at (110,\base) {$+$};
    \draw[mc flow,draw=NNOutput] (kernel\o3-out)
      to[out=0,in=145] (sum\o.145);
    \draw[mc flow,draw=NNHidden] (kernel\o2-out)--(sum\o.west);
    \draw[mc flow,draw=NNInput] (kernel\o1-out)
      to[out=0,in=215] (sum\o.215);
  }
  % 输入从每个切片的左侧进入，切片的结果只从右侧汇出。
  \foreach \c/\col in {1/NNInput,2/NNHidden,3/NNOutput}{
    \draw[mc flow,draw=\col] (input\c)
      to[out=0,in=180,looseness=1.05] (kernel1\c-in);
    \draw[mc flow,draw=\col] (input\c)
      to[out=0,in=180,looseness=1.05] (kernel2\c-in);
  }
  \node[nn heading,anchor=south] at (78.5,46) {核1};
  \node[nn heading,anchor=north] at (78.5,-46) {核2};
  \node (bias1) at (110,44) {$b_1$};
  \node (bias2) at (110,-44) {$b_2$};
  \draw[mc flow,draw=NNHidden] (bias1.south)--(sum1.north);
  \draw[mc flow,draw=NNHidden] (bias2.north)--(sum2.south);
  % 两个输出通道同样错位叠放，曲线各自到达对应响应。
  \foreach \o/\xx/\yy/\col in {1/145/1/NNHidden,2/136/-8/NNOutput}{
    \draw[draw=\col,fill=\col!7,line width=1.1pt]
      (\xx,\yy) rectangle ({\xx+18},{\yy+18});
    \foreach \g in {1,2}{
      \draw[draw=\col!45,line width=.5pt]
        ({\xx+6*\g},\yy)--({\xx+6*\g},{\yy+18})
        (\xx,{\yy+6*\g})--({\xx+18},{\yy+6*\g});
    }
    \draw[draw=\col,fill=\col!34,line width=1.05pt]
      ({\xx+12},{\yy+12}) rectangle ({\xx+18},{\yy+18});
  }
  \draw[mc flow,draw=NNHidden] (sum1.east)
    .. controls (131,26) and (135,16) .. (157,16);
  \draw[mc flow,draw=NNOutput] (sum2.east)
    .. controls (133,-26) and (127,7) .. (148,7);
  \node[mc channel] at (158.5,5) {$o=1$};
  \node[mc channel] at (144,-5) {$o=2$};
  \node[nn heading,anchor=north] at (149.5,-12) {$\bm Z$};
\end{tikzpicture}
